% !TEX root = ../../tree_notes.tex
% =====================================================================
%  Chapter 10 opener -- "Chapter at a Glance", three reading paths + the
%  ladder staircase, the two running examples (the array R = 5,3,8,1,4,2,7
%  and the stone pile S = 2,7,4,1,8,1), the "who owns what" comparison
%  (heap vs BST vs Cartesian tree vs segment tree), roadmap, recognition
%  guide, constraint->tool table, notation.
%  NO numbered boxes here. Exports fig:c07-layout and fig:c07-heapR.
% =====================================================================

\begin{chaptertheorybox}{Chapter at a Glance}
Chapter~9 (\chref{c26}) \emph{read} a heap-order straight out of a static array -- the Cartesian tree,
built once by a monotonic stack and never touched again. This chapter is the moving counterpart. A
\term{binary heap} is a complete binary tree stored in a plain array that you \emph{edit}: you push a
value, you pop the extreme, you sift things back into place -- and the whole time you know, in one array
cell, what the largest (or smallest) element is. The promise is deliberately weak: the root beats its
children, but \emph{siblings are unordered}. That one weak promise -- ``the top is the extreme'' -- is
worth an enormous amount, because it is all greedy extraction, top-$k$ selection, the two-heap running
median, and multiway merge ever ask for:
\begin{itemize}[nosep]
  \item \textbf{An array in disguise.} A complete tree needs no pointers: element $i$ has parent
        $\lfloor i/2\rfloor$ and children $2i$ and $2i+1$. The heap is ``a handful of array indices and
        one \code{while} loop.'' \hfill \lv{2}--\lv{3}
  \item \textbf{Edit, don't reorder.} \code{push} plants a leaf and sifts \emph{up}; \code{pop} moves the
        root out, drops the last element into it, and sifts \emph{down}. Both are $\bigO{\log n}$ because
        the tree stays \emph{short by shape}, not by rotations. \hfill \lv{3}
  \item \textbf{Build is $\bigO{n}$, not $\bigO{n\log n}$.} Heapifying bottom-up (Floyd) pays a geometric
        series -- most nodes sit near the leaves and move barely at all. \hfill \lv{4}
\end{itemize}

\textbf{You will learn to}
\begin{itemize}[nosep]
  \item store a heap in an array and write sift-up / sift-down push and pop cold; \hfill\lv{3}
  \item build a heap in $\bigO{n}$ and explain why the naive $\bigO{n\log n}$ is not optimal; \hfill\lv{4}
  \item drive \code{std::priority\_queue} -- min-heap via \code{greater}, custom comparators, the
        \code{pair} trick -- without hand-rolling what the library already does; \hfill\lv{3}
  \item keep top-$k$ / $k$-th answers in a size-$k$ heap, and a running median in \emph{two} heaps; \hfill\lv{4}
  \item merge heaps cheaply with a leftist/skew structure and know when binomial/Fibonacci heaps matter. \hfill\lv{4}
\end{itemize}

\textbf{Prerequisites.} Arrays and index arithmetic; the complete-binary-tree shape previewed in
Chapter~4 (\chref{c04}); tree vocabulary (root, height, leaf) from Chapter~1. No recursion is required
beyond the height bound, and no search-tree ordering is assumed -- indeed \secref{c07-heap} opens by
contrasting the heap's ``no sibling order'' against the BST inorder of Chapter~7 (\chref{c05}). The
graph use of a priority queue (Dijkstra/Prim) is \emph{previewed, not taught}.
\end{chaptertheorybox}

% ---------------- Reading paths + ladder position ----------------
\begin{chaptertheorybox}{How to Read This Chapter (three paths)}
Level badges $\lv1$--$\lv6$: $\lv2$ is a bare array-layout fact, $\lv3$ is a solid LC-Easy/Medium heap
problem you solve with \code{priority\_queue} out of the box, $\lv4$ is where you must \emph{design} the
heap (a size-$k$ heap, two heaps, a merge), and $\lv5$ is the Summit, where the whole difficulty is
finding the right greedy invariant.

\textbf{Fast path ($\lv3$, ``I just want a priority queue'').} Read \secref{c07-heap} and the
\code{std::priority\_queue} corner of \secref{c07-build}. Do the solved Last Stone Weight (LC 1046) and
the top-$k$ lab (LC 347), then the practices. Skip the Floyd proof, the mergeable-heap section, and the
treap-adjacent material. You leave able to cold-solve $k$-th-largest, top-$k$-frequent, and any
``greedily grab the extreme'' problem.

\textbf{Structural path ($\lv3$--$\lv4$).} Add \secref{c07-build}'s $\bigO{n}$-build proof and
\secref{c07-trees} on mergeable heaps -- why a binary heap \emph{cannot} merge cheaply and what a
leftist/skew heap does about it. Order: 7.1 $\to$ 7.2 $\to$ 7.3.

\textbf{Full path.} Everything in order, finishing at \secref{c07-classics} (multiway merge, the two-heap
running median, best-first $k$-pairs) and the Summit (CF 721D), which is the purest ``extract-min and
reinsert'' loop: a min-heap keyed by absolute value, driven $k$ times to push a product as negative as
possible.

The forward rope is Chapter~11 (\chref{c08}) for tries, and the graph chapters for the priority queue used
as a \emph{scheduling} engine; the backward rope is Chapter~9 (\chref{c26}) for the read-only heap order
you do \emph{not} maintain.
\end{chaptertheorybox}

\begin{figure}[H]
\centering
\begin{tikzpicture}[diagram, x=9mm, y=7mm]
  \foreach \st/\lab in {1/lv1,2/lv2,3/lv3,4/lv4,5/lv5,6/lv6} {
    \fill[\lab!22, draw=\lab, line width=0.8pt] (\st-1,0) rectangle (\st,\st*0.52);
    \node[font=\scriptsize\bfseries, text=\lab] at (\st-0.5,\st*0.52-0.42) {L\st};}
  \node[annot, anchor=south west, font=\scriptsize]
    at (0.05,0.62) {a tree that is just an array};
  \node[annot, anchor=south west, font=\scriptsize]
    at (1.05,1.14) {7.1 layout; push/pop};
  \node[annot, anchor=south west, font=\scriptsize]
    at (2.05,1.66) {7.2 Floyd build; priority\_queue};
  \node[annot, anchor=south west, font=\scriptsize]
    at (3.05,2.18) {7.2 top-k; 7.3 mergeable};
  \node[annot, anchor=south west, font=\scriptsize, text=lv5]
    at (4.05,2.70) {7.4 multiway merge; two-heap median};
  \node[annot, anchor=south west, font=\scriptsize, text=lv5]
    at (5.05,3.22) {summit (CF 721D)};
\end{tikzpicture}
\caption*{Chapter~10 on the six-step ladder: the array layout and sift are $\lv2$--$\lv3$; the $\bigO{n}$
build and the library priority queue sit at $\lv3$--$\lv4$; designing a size-$k$ or two-heap structure
climbs to $\lv4$; multiway merge and the running median reach $\lv5$, where the Summit (CF 721D
\emph{Maxim and Array}) turns the whole chapter into one extract-min-and-reinsert loop. Fast path $=$
steps 2--4; full path runs to 6.}
\end{figure}

% ---------------- The running examples ----------------
\subsection*{The Running Examples: One Array, One Stone Pile}

Two small inputs, reused in every section so the ideas always land on something concrete.

\textbf{Array $R$} $= (5, 3, 8, 1, 4, 2, 7)$, indices $1,\dots,7$ -- seven values, a perfect complete tree
of depth~$2$. Laid out into an array it is just numbers; read with the index rule it is a tree; heapify
it and you get the min-heap $(1, 3, 2, 5, 4, 8, 7)$ whose root $1$ is the minimum, but whose two children
$3$ and $2$ are \emph{out of any left-to-right order} -- the exact promise a heap makes. $R$ drives the
layout and sift traces (\secref{c07-heap}) and the Floyd build (\secref{c07-build}).

\textbf{Stone pile $S$} $= (2, 7, 4, 1, 8, 1)$ -- the official LeetCode~1046 input. You repeatedly smash
the two heaviest stones, leaving their difference; the answer is the last stone's weight ($\boxed{1}$ for
$S$). ``Always take the two largest'' is precisely the \code{pop}-the-extreme operation, so $S$ becomes
the solved lab of \secref{c07-heap} and reappears as the template for every greedy-extraction problem.

\begin{figure}[H]
\centering
\begin{tikzpicture}[diagram, x=13mm, y=9mm, level distance=13mm, sibling distance=22mm]
  % index layout of the complete tree
  \node[boxnode, fill=primaryblue!12, draw=primaryblue] {$i{=}1$}
    child { node[boxnode, fill=white, draw=darkgray] {$i{=}2$}
      child { node[boxnode, fill=white, draw=darkgray] {$i{=}4$} }
      child { node[boxnode, fill=white, draw=darkgray] {$i{=}5$} } }
    child { node[boxnode, fill=white, draw=darkgray] {$i{=}3$}
      child { node[boxnode, fill=white, draw=darkgray] {$i{=}6$} }
      child { node[boxnode, fill=white, draw=darkgray] {$i{=}7$} } };
  \node[annot, align=left, font=\scriptsize] at (-3.4,-2.2)
    {parent of $i$ $=\lfloor i/2\rfloor$;\\ children $2i,\;2i{+}1$;\\ no pointers, just indices};
  \node[annot, align=left, font=\scriptsize] at (3.6,-2.2)
    {the array $R$ written\\ into the complete-tree\\ slots $1\ldots 7$};
\end{tikzpicture}
\caption{The index layout a heap exploits. A complete binary tree of $n$ nodes fills positions
$1..n$ with no gaps, so a child is reachable by $2i,2i+1$ and a parent by $\lfloor i/2\rfloor$ -- the
whole structure is one array and integer arithmetic (\secref{c07-heap}). This shape is fixed: a heap is
never rebalanced by rotations, only \emph{refilled} at the last slot.}\label{fig:c07-layout}
\end{figure}

\begin{figure}[H]
\centering
\begin{tikzpicture}[diagram, x=13mm, y=9mm, level distance=13mm, sibling distance=22mm]
  \node[troot] {$1$}
    child { node[tnode] {$3$}
      child { node[tleaf] {$5$} }
      child { node[tleaf] {$4$} } }
    child { node[tnode] {$2$}
      child { node[tleaf] {$8$} }
      child { node[tleaf] {$7$} } };
  \node[annot, align=left, font=\scriptsize] at (-4.4,-0.5)
    {\textbf{Min-heap on $R$:}\\ array $(1,3,2,5,4,8,7)$\\ parent $\le$ both children};
  \node[annot, align=left, font=\scriptsize] at (4.2,-2.0)
    {\textbf{Siblings unordered:}\\ left child $3 >$ right child $2$\\ -- no inorder, no BST};
\end{tikzpicture}
\caption{Array $R=(5,3,8,1,4,2,7)$ after Floyd build, as the min-heap $(1,3,2,5,4,8,7)$. The root $1$ is
the minimum and every parent is $\le$ its children, but the two children of the root are $3$ and $2$ --
\emph{out of order}. A heap guarantees only ``parent beats children''; the left-to-right reading is
meaningless, which is exactly what separates it from the BST of Chapter~7 and the Cartesian tree of
Chapter~9 (\secref{c07-heap}).}\label{fig:c07-heapR}
\end{figure}

\medskip
\noindent Four structures answer superficially similar ``find the extreme'' questions; knowing which owns
what stops you reaching for the wrong engine:
\begin{center}
\renewcommand{\arraystretch}{1.25}
\setlength{\tabcolsep}{5pt}
\begin{tabular}{@{}p{0.185\linewidth} p{0.42\linewidth} p{0.31\linewidth}@{}}
\toprule
\textbf{Structure} & \textbf{What it is for} & \textbf{Where it lives} \\
\midrule
Binary heap / \code{priority\_queue} & the \emph{extreme} of a changing multiset, plus repeated
  ``extract-min/max'' for greedy or top-$k$ & this chapter (\secref{c07-heap}, \secref{c07-build}) \\
Cartesian tree (read-only heap order) & the min/max structure \emph{already in a static array}, read
  out by one stack, never edited & Chapter~9 (\chref{c26}) \\
BST / balanced tree & \emph{ordered by key}: predecessor, successor, any rank -- a total order, not just
  an extreme & Chapters~7--8 (\chref{c05}, \chref{c06}) \\
Segment / sparse table & \emph{arbitrary range} aggregates over an array, answering any $[l,r]$, not only
  the whole & Chapters~18--19 (\chref{c14}, \chref{c13}) \\
\bottomrule
\end{tabular}
\end{center}

% ---------------- Roadmap diagram ----------------
\begin{figure}[H]
\centering
\begin{tikzpicture}[diagram, node distance=3.5mm,
    top/.style={boxnode, fill=primaryblue!10, draw=primaryblue, text width=21mm,
                minimum height=11mm, font=\small},
    res/.style={annot, text width=21mm, anchor=north, font=\scriptsize}]
  \node[top] (t1) {\textbf{7.1} The Binary Heap};
  \node[top, right=of t1] (t2) {\textbf{7.2} Build + Use};
  \node[top, right=of t2] (t3) {\textbf{7.3} Mergeable};
  \node[top, right=of t3] (t4) {\textbf{7.4} Classics};
  \foreach \a/\b in {t1/t2, t2/t3, t3/t4} \draw[flow] (\a) -- (\b);
  \node[res] at ([yshift=-2mm]t1.south) {array layout; sift up/down};
  \node[res] at ([yshift=-2mm]t2.south) {Floyd $\bigO{n}$; \code{priority\_queue}};
  \node[res] at ([yshift=-2mm]t3.south) {leftist/skew merge; costs};
  \node[res] at ([yshift=-2mm]t4.south) {merge-$k$; two-heap median};
\end{tikzpicture}
\caption*{The four topics are one structure seen at four depths. 7.1 \emph{lays} the tree in an array and
sifts; 7.2 builds it in $\bigO{n}$ and hands you the library queue plus the top-$k$ pattern; 7.3 asks
``can I merge two heaps fast?'' and answers with leftist/skew trees (binomial/Fibonacci as ideas); 7.4
spends the tool on the classics -- multiway merge, running median, best-first $k$-pairs -- that make the
Summit's loop feel natural.}
\end{figure}

% ---------------- Topic table ----------------
\begin{center}
\small
\renewcommand{\arraystretch}{1.25}
\setlength{\tabcolsep}{4pt}
\begin{tabular}{@{}p{0.235\linewidth} p{0.365\linewidth} c p{0.23\linewidth}@{}}
\toprule
\textbf{Topic} & \textbf{Key results} & \textbf{Levels} & \textbf{Solved problem} \\
\midrule
\secref{c07-heap}\ \ The Binary Heap
  & array layout, height $\lfloor\lg n\rfloor$, sift-up/down, push/pop $\bigO{\log n}$ & $\lv2$--\lv{3}
  & \probref{c07-stone} \\
\addlinespace
\secref{c07-build}\ \ Floyd Build + the Queue
  & $\bigO{n}$ heapify (geometric series); \code{priority\_queue}; size-$k$ top-$k$ & $\lv3$--\lv{4}
  & \probref{c07-topk} \\
\addlinespace
\secref{c07-trees}\ \ Mergeable Heaps
  & leftist npl, short right spine; merge cost table; Huffman-greedy & $\lv4$
  & \probref{c07-connect} \\
\addlinespace
\secref{c07-classics}\ \ Classic Heap Problems
  & multiway merge, two-heap median, best-first $k$-pairs & $\lv4$--\lv{5}
  & \probref{c07-mergek} \\
\bottomrule
\end{tabular}
\end{center}

% ---------------- Problem recognition guide ----------------
\subsection*{Problem Recognition Guide --- When You See \emph{this}, Think \emph{that}}

Heap problems are the ones where you repeatedly need the \emph{current} extreme of a set that changes.
If the phrasing mentions ``largest/smallest remaining'', ``top $k$'', ``merge the two smallest'', or
``schedule by priority'', that is the tell; every row points at the section that builds the tool.

\begin{center}
\renewcommand{\arraystretch}{1.3}
\setlength{\tabcolsep}{4pt}
\begin{tabular}{@{}p{0.44\linewidth} p{0.34\linewidth} l@{}}
\toprule
\textbf{When the problem says\ldots} & \textbf{\textcolor{deepblue}{\textbf{Think\ldots}}} & \textbf{Section} \\
\midrule
"the largest / smallest so far", "the $k$-th largest", "keep the top $k$"
  & a (size-$k$) heap; push then pop to hold the frontier & \secref{c07-build} \\
\addlinespace
"always take the two extremes and combine them" (smash stones, connect ropes, Huffman)
  & greedy extraction: pop, merge, push back & \secref{c07-heap}, \secref{c07-trees} \\
\addlinespace
"median of a stream" / "balance two halves"
  & two heaps (a max-heap low, min-heap high) & \secref{c07-classics} \\
\addlinespace
"merge $k$ sorted sequences", "smallest range hitting all lists"
  & a min-heap of one head per list (multiway merge) & \secref{c07-classics} \\
\addlinespace
"k-th smallest pair sum / in a sorted matrix"
  & best-first expansion of a min-heap & \secref{c07-classics} \\
\addlinespace
"must merge two priority queues fast"
  & a mergeable heap (leftist/skew), not a binary one & \secref{c07-trees} \\
\bottomrule
\end{tabular}
\end{center}

\begin{center}
\renewcommand{\arraystretch}{1.25}
\setlength{\tabcolsep}{4pt}
\begin{tabular}{@{}p{0.29\linewidth} p{0.42\linewidth} p{0.20\linewidth}@{}}
\toprule
\textbf{Constraint you see} & \textbf{What it forces} & \textbf{Chapter-10 tool} \\
\midrule
$10^{5}$--$10^{6}$ pushes/pops, ``extreme at each step'' & $\bigO{\log n}$ per op -- a sorted array would & \code{priority\_queue} \\
$k \ll n$, "the top $k$ / $k$-th largest" & keep only $k$ in a heap, don't sort all & size-$k$ heap \\
repeatedly combine the two smallest & greedy with a min-heap, $\bigO{n\log n}$ & \probref{c07-connect} \\
merge $k$ lists totalling $N$ nodes & $\bigO{N\log k}$, a heap of $k$ heads & \probref{c07-mergek} \\
streaming median, only the middle needed & two heaps, sizes within $1$ & \secref{c07-classics} \\
\bottomrule
\end{tabular}
\end{center}

\begin{notebox}[Not this chapter?]
If you need the extreme of a \emph{static} array and never edit it, that is the Cartesian tree / monotonic
stack of Chapter~9 (\chref{c26}) -- no \code{priority\_queue} needed. If you need \emph{any} rank
($k$-th smallest \emph{and} largest, predecessor, successor), a heap is too weak -- you need an ordered
tree (Chapters~7--8) or a Fenwick/segment tree over compressed keys (\chref{c14}, \chref{c13}). If you
need the extreme of every \emph{subarray} $[l,r]$, that is range-min, not a global heap (\chref{c26},
\chref{c13}). A heap owns exactly one thing and owns it perfectly: \emph{the current extreme of a
multiset you push into and pop from}.
\end{notebox}

% ---------------- Notation ----------------
\subsection*{Notation Used in This Chapter}

\begin{center}
\renewcommand{\arraystretch}{1.2}
\setlength{\tabcolsep}{4pt}
\begin{tabular}{@{}p{0.20\linewidth} p{0.27\linewidth}@{\quad}p{0.17\linewidth} p{0.27\linewidth}@{}}
\toprule
\textbf{Symbol} & \textbf{Meaning} & \textbf{Symbol} & \textbf{Meaning} \\
\midrule
$R$                 & the running array $(5,3,8,1,4,2,7)$ (\figref{c07-heapR}) & $S$ & the stone pile $(2,7,4,1,8,1)$ (LC 1046) \\
$a[i]$              & value at index $i$; heaps here are $1$-based & $\mathit{pq}$ & a \code{priority\_queue} (a heap) \\
$\lfloor i/2\rfloor$ & parent of node $i$; children are $2i$ and $2i+1$ & $\mathit{size}$ & number of elements currently in the heap \\
$\height(v)$        & height of the subtree at $v$ & $\mathit{npl}$ & null-path length of a leftist node (\secref{c07-trees}) \\
$\lg n$             & $\log_2 n$; a heap's height is $\lfloor\lg n\rfloor$ & $\bigstar$ & marks the Summit \\
$\min/\max$ heap    & root is the smallest / largest element; \code{parent}\,$\le$\,\code{child} & & \\
\bottomrule
\end{tabular}
\end{center}
